% Online-only extension; implementation and protocol in ../orin_scheduling/.
\begingroup
\color{blue}
\subsubsection{Feature Caching}
\label{subsec:feature-prefetch}

An idle gap between foreground requests may be too short for a complete training batch, but long enough for its frozen forward computation.
We cache the output just before the first trainable layer to schedule these two parts separately.
The frozen forward computation can finish in one gap, and training from the first trainable layer can resume in a later gap using the stored features.
We let each admitted operation finish without interrupting it.

The cache is a temporary buffer within one epoch.
It starts empty, and each feature is consumed once and then released.
The next epoch recomputes features even for the same sample and proxy.
Thus, caching does not remove forward computations through reuse across epochs; it allows their results to wait until the remaining training work can run.

Caching also allows us to prepare features for later batches.
For a fixed input and proxy, the frozen output does not depend on the current trainable parameters.
We can therefore reorder feature computations while keeping the work from the first trainable layer onward in the original batch order.
The original sampler sets the batch plan at the start of each epoch, and subsequent loss calculations, update decisions, and sampler feedback use the current parameters and loss history.

If the next batch's features are ready, its remaining training work takes priority; if that work cannot fit, we wait.
Otherwise, we first try to compute its missing features.
If they cannot fit, we look ahead and choose the longest fitting feature computation for another proxy.
We use a fixed-capacity buffer and reserve space for the earliest missing batch.
We admit each operation only when its estimated runtime fits before the next foreground request.
The cache stores features computed ahead of time, while Reorder chooses which ones to prepare.
Its scheduling benefit comes from using gaps that would otherwise be left idle.
Grouping batches that share a proxy can additionally improve forward efficiency, a benefit that does not require foreground inference.
Offline proxy generation and policy training stay unchanged.

\par
\endgroup
